\begin{helper}
Fix an activation set $A$ and a vertex $v$ in the sampling space of
\noderef{RandomIndependentSetSampling}.  Let $\nu$ be any measure satisfying
the rectangular priority clause from that definition.  If $0\le a\le 1$, then
the $\nu$-measure of the rectangle in which the activation set is exactly
$A$, every priority coordinate lies in $[0,1]$, and each activated neighbor
of $v$ has priority at most $a$ is
\[
\nu[\omega_1=A]\,
a^{|\{u\in V(G):u\in A\text{ and }u\sim_G v\}|}.
\]
\end{helper}
\begin{proof}
Define a threshold function $t$ on vertices by
\[
t(u)=
\begin{cases}
a, & u\in A\text{ and }u\sim_G v,\\
1, & \text{otherwise}.
\end{cases}
\]
Because $0\le a\le 1$, every value of $t$ lies in $[0,1]$.  Applying the
rectangular priority clause of \noderef{RandomIndependentSetSampling} to
this threshold function gives
\[
\nu\{\omega:\omega_1=A\text{ and }0\le \omega_2(u)\le t(u)
\text{ for every }u\}
=
\nu[\omega_1=A]\prod_{u\in V(G)}t(u).
\]
The event on the left is exactly the event in the statement: the coordinates
with $u\in A$ and $u\sim_G v$ are bounded above by $a$, while all other
coordinates are only bounded above by $1$.

It remains only to simplify the finite product.  Each factor is $a$ precisely
for vertices in the finite set
\[
\{u\in V(G):u\in A\text{ and }u\sim_G v\},
\]
and each remaining factor is $1$.  Hence the product is
$a^{|\{u\in V(G):u\in A\text{ and }u\sim_G v\}|}$, which is the claimed
formula.
\end{proof}
